\begin{tikzpicture}[figure picture,x=1mm,y=1mm,
  causalnode/.style={circle,draw=FigureStroke,fill=none,line width=1pt,
    minimum size=7mm,inner sep=0pt,font=\figurefont\mdseries\footnotesize},
  conditionednode/.style={causalnode,draw=LancetGreen,
    double=white,double distance=0.7pt},
  causaledge/.style={-{Stealth[length=2.2mm,width=1.4mm]},
    draw=FigureStroke,line width=1.1pt},
  paneltitle/.style={font=\figurefont\mdseries\footnotesize,text=FigureInk},
  panellabel/.style={font=\figurefont\mdseries\small,text=FigureInk},
  note/.style={font=\figurefont\mdseries\footnotesize,
    text=FigureMuted,align=center,fill=none}]
  \path[use as bounding box] (0,-10) rectangle (169,38);
  \draw[FigureGrid,line width=0.6pt] (56,-8)--(56,36);
  \draw[FigureGrid,line width=0.6pt] (112,-8)--(112,36);

  \node[panellabel,anchor=west] at (1,33) {(a)};
  \node[paneltitle] at (28,33) {共同原因};
  \node[conditionednode] (fork-x) at (28,22) {\(X\)};
  \node[causalnode] (fork-t) at (13,11) {\(T\)};
  \node[causalnode] (fork-y) at (43,11) {\(Y\)};
  \draw[causaledge] (fork-x)--(fork-t);
  \draw[causaledge] (fork-x)--(fork-y);
  \node[note] at (28,-3) {给定\(X\)时关闭路径};

  \node[panellabel,anchor=west] at (57,33) {(b)};
  \node[paneltitle] at (84,33) {中介链};
  \node[causalnode] (chain-t) at (67,16) {\(T\)};
  \node[conditionednode] (chain-m) at (84,16) {\(M\)};
  \node[causalnode] (chain-y) at (101,16) {\(Y\)};
  \draw[causaledge] (chain-t)--(chain-m);
  \draw[causaledge] (chain-m)--(chain-y);
  \node[note] at (84,-3) {给定\(M\)时关闭路径};

  \node[panellabel,anchor=west] at (113,33) {(c)};
  \node[paneltitle] at (140,33) {碰撞点};
  \node[causalnode] (collider-t) at (125,22) {\(T\)};
  \node[causalnode] (collider-y) at (155,22) {\(Y\)};
  \node[conditionednode] (collider-c) at (140,11) {\(C\)};
  \draw[causaledge] (collider-t)--(collider-c);
  \draw[causaledge] (collider-y)--(collider-c);
  \node[note] at (140,-3) {不给定时关闭\\给定\(C\)时打开};
\end{tikzpicture}
